% document formatting
\documentclass[10pt]{article}
\usepackage[utf8]{inputenc}
\usepackage[left=1in,right=1in,top=1in,bottom=1in]{geometry}
\usepackage[T1]{fontenc}
\usepackage{xcolor}

% math symbols, etc.
\usepackage{amsmath, amsfonts, amssymb, amsthm}

% lists
\usepackage{enumerate}

% images
\usepackage{graphicx} % for images
\usepackage{tikz}

% code blocks
% \usepackage{minted, listings} 
\usepackage[colorlinks=true, urlcolor=blue]{hyperref}


% verbatim greek
\usepackage{alphabeta}

\graphicspath{{./assets/images/Week 7}}

\title{02-719 Week 7 \\ \large{Genomics and Epigenetics of the Brain}}
 
\author{Aidan Jan}

\date{\today}

\begin{document}
\maketitle
\section*{t-SNE}
First, let's consider a PCA of a cell dataset.\\\\
Recall PCA is a dimensionality reduction technique.
\begin{itemize}
	\item Linear
	\item Flavor of Matrix Factorization
	\item Maximizes variance explained by components
\end{itemize}
\begin{center} 
	\includegraphics*[width=0.6\textwidth]{W7_1.png} 
\end{center}
Here, each point is a cell, and each color is a cell type in the retina.
\begin{itemize}
	\item Why doesn't this look good?
	\item Sometimes, the data isn't linearly separable!
\end{itemize}
Now, let's try the nonlinear method, t-SNE:
\begin{center} 
	\includegraphics*[width=0.6\textwidth]{W7_2.png} 
\end{center}
\begin{itemize}
	\item This is much better!
	\item Unlike PCA, t-SNE is able to capture nonlinear dynamics.  (These are very prevalent in snRNA-req data).
	\item Each PC in snRNA-seq data usually captures $<1\%$ of the total variance
	\item PCA focuses on global variance so the first few components may not include information on more granular structure
\end{itemize}

\subsection*{How can similarity be defined?}
It is usually defined based on the PCA metagene matrix we made
\begin{itemize}
	\item Cosine distance/similarity
	\item Euclidean distance can be turned into similarity if taken through a gaussian kernel
	\item Number of shared nearest neighbors can define similarity
\end{itemize}
The similarity matrix can be viewed as a graph
\begin{itemize}
	\item Nodes are cells
	\item Edges are the similarity between 2 cells.
\end{itemize}
\begin{center} 
	\includegraphics*[width=0.6\textwidth]{W7_3.png} 
\end{center}

\subsection*{TNSE-T distributed stochastic neighbor embedding}
\begin{itemize}
	\item Tries to replicate the similarity found in the similarity matrix in two dimensions
	\item Actually a probabilistic approach
	\begin{itemize}
	    \item Given a distance between two cells it creates a probability that they should be distance $X$ away from each other in the 2d space
    \end{itemize}
    \item This probability is based on the t-distribution with defined parameter perplexity
\end{itemize}

\subsection*{TNSE - How perplexity can change the visualization}
\begin{center} 
	\includegraphics*[width=0.8\textwidth]{W7_4.png} 
\end{center}

\section*{UMAP}
UMAP is an alternative to TNSE.
\begin{itemize}
	\item Instead of the probabilistic distance between cells this focuses on the topology or graph structure of that similarity matrix and tries to maintain it in the 2d projection
	\item Two parameters:
	\begin{itemize}
	    \item \texttt{min\_dist}: the minimum distance considered meaningful
	    \item \texttt{N\_neighbors}: how far out in the graph does structure matter?
    \end{itemize}
\end{itemize}
\subsection*{UMAP on a 3D Mammoth}
\begin{center} 
	\includegraphics*[width=0.9\textwidth]{W7_5.png} 
\end{center}

\subsection*{Comparison of Dimensionality Reduction Methods}
\begin{center}
\begin{tabular}{c|c}
{\underline{\textbf{PCA}}}                                                                         & {\underline{\textbf{UMAP / TSNE}}}                                                                                              \\ \hline
Linear                                                                                             & Non-linear                                                                                                              \\ \hline
\begin{tabular}[c]{@{}c@{}}Dimensions are explainable\\ by gene embeddings\end{tabular}            & \begin{tabular}[c]{@{}c@{}}Dimensions have no\\ intrinsic meaning\end{tabular}                                          \\ \hline
\begin{tabular}[c]{@{}c@{}}Deterministic (will get the \\ same dimensions every time)\end{tabular} & Non-deterministic                                                                                                       \\ \hline
\begin{tabular}[c]{@{}c@{}}Captures small amount of \\ variance per dimension\end{tabular}         & \begin{tabular}[c]{@{}c@{}}2-3 dimensions can give a \\ good look at the variability\\ across your dataset\end{tabular}
\end{tabular}
\end{center}
Note that \textbf{UMAP and t-SNE are \textit{approximations} of the similarity matrix}.
\begin{itemize}
	\item They should be used for visualization, data exploration, and hypothesis generation
	\item They should not be used to draw statistical conclusions about your data
\end{itemize}

\section*{Clustering}
\subsection*{Louvain Community Detection}
\begin{itemize}
	\item Community detection is a type of clustering performed on a similarity graph
	\item Find communities (clusters) of tightly knit nodes
	\item Optimizes modularity
\end{itemize}

\subsection*{Louvain Algorithm Steps}
\begin{itemize}
	\item The Louvain algorithm consists of 2 repeating steps:
	\begin{itemize}
	    \item Step 1 moves nodes to the optimal community
	    \item Step 2 aggregates all nodes in a community
    \end{itemize}
    \item Each repetition of these steps is called a level
    \item The steps are repeated until convergence (nothing changes between levels)
\end{itemize}
\begin{center} 
	\includegraphics*[width=0.6\textwidth]{W7_6.png} 
\end{center}

\subsubsection*{Step 0: Initialization}
Each vertex $v_i$ has its own community $c_i$.\\\\

\subsubsection*{Step 1}
Step 1 can be divided into passes.  In each pass, the following algorithm is used:
\begin{itemize}
	\item For each vertex $v_i$,
	\begin{itemize}
	    \item For each community $c_j$
	    \begin{itemize}
	        \item If adding vertex $v_i$ to community $c_j$ increases the modularity of the similarity graph, add $v$ to $c$.
        \end{itemize}
    \end{itemize}
\end{itemize}

\subsubsection*{Step 2}
Step 2 is the merge step - all the vertices in each community are merged into a single vertex.  The edges between communities are the sum of edges between the vertices in the communities.  After this, repeat the process until convergence.

\subsubsection*{Adding a Resolution Parameter}
\begin{itemize}
	\item You will likely want to cluster your single cell data set at different resolutions.
	\item To address this we can add a resolution parameter to weight $p_{ij}$
	\item The higher the resolution the more fine grained the clusters will become
	\[Q = \frac{1}{2m} \sum_{i, j} \left[A_{i, j} - r\frac{k_i k_j}{2m}\right] \delta(c_i, c_j)\]
\end{itemize}

\subsection*{Effects of the Resolution Parameter}
\begin{center} 
	\includegraphics*[width=0.8\textwidth]{W7_7.png} 
\end{center}

\subsection*{Marker Gene Plots}
\begin{center} 
	\includegraphics*[width=0.9\textwidth]{W7_8.png} 
\end{center}

\subsection*{Rbfox 3 on a Violin Plot Mouse Motor Cortex Cells}
Let's use an example: Rbfox3 is a marker gene for neurons.
\begin{center} 
	\includegraphics*[width=0.9\textwidth]{W7_9.png} 
\end{center}

\subsection*{Rbfox3 on a UMAP of Mouse Motor Cortex Cells}
\begin{center} 
	\includegraphics*[width=\textwidth]{W7_10.png} 
\end{center}

\subsection*{Marker Genes are Relative}
\begin{itemize}
	\item A marker gene allows a cell type to be identifiable compared to every other cell type in the population you are analyzing
	\item If you have the same cell type with a different population in your analysis, you may get different marker genes!
\end{itemize}

\subsection*{Differential Expression Tests}
\begin{itemize}
	\item A differential expression statistical test is used to see which genes differ in their expression distribution between 2 populations
	\item It has the following null hypothesis for each gene $X$: the expression of gene $X$ was drawn from the same distribution for both population $A$ and $B$.
	\item To perform these tests we must first normalize the expression values
\end{itemize}

\subsection*{Wilcoxon-Rank Sum Test}
\begin{itemize}
	\item The T-test is a \textit{parametric} test, which means it makes assumptions about the types of distributions that the statistic is calculated from.  Namely, the t-test assumption is the difference between expression distributions is gaussian.
	\item The Wilcoxon-rank sum test makes no such assumptions.  It is \textit{non-parametric}.
	\item Instead of using mean or variance measurement it uses a difference in ranks for its statistic.
\end{itemize}

\subsection*{More Complex Experimental Designs}
\begin{itemize}
	\item For these marker genes we only looked at the difference between two groups.
	\item Why may this be too simplistic?
	\begin{itemize}
	    \item We might want to take into account other factors such as: sequencing batch, sex of sample, age/time point at which the sample was taken, multiple treatment conditions, disease progression, etc.
    \end{itemize}
\end{itemize}

\subsection*{Linear Models}
Linear models allow you to take into account multiple factors when calculating differential expression.
\begin{itemize}
	\item Rbfox3 expression $= \beta_0 + \beta_1 + \beta_2 + \beta_3$, where $\beta_1 :=$ batch, $\beta_2 :=$ sex, $\beta_3 :=$ is neuron.
\end{itemize}
In this case our statistic will be $\beta_3$ and our $p$-value will be calculated based off of it.  There are many packages to use linear models for RNA-seq data, for single cell recommend limma-trend.


\end{document}